% Part: methods
% Chapter: proofs
% Section: introduction

\documentclass[../../../include/open-logic-section]{subfiles}

\begin{document}

\olfileid{mth}{prf}{int}

\olsection{ভূমিকা}

যুক্তিবিদ্যার প্রাথমিক পাঠের অভিজ্ঞতা থেকে তুমি হয়তো
!!a{derivation} পদ্ধতিতে স্বচ্ছন্দ—সম্ভবত স্বাভাবিক নিষ্পাদন বা
ফিচ-শৈলীর !!{derivation} পদ্ধতিতে, কিংবা প্রমাণ-বৃক্ষের কোনো
পদ্ধতিতে। এসব পদ্ধতিতে প্রমাণ করার কথাও নিশ্চয় মনে আছে:
!!a{formula} প্রমাণ করা অথবা প্রদত্ত যুক্তি বৈধ দেখানো।
কাঙ্ক্ষিত সিদ্ধান্তে পৌঁছানো পর্যন্ত পদ্ধতির বিধি প্রয়োগ করেই
তা করতে। যুক্তিবিদ্যা \emph{সম্বন্ধে} বিচার করতেও আমরা প্রমাণ দিই,
কিন্তু অধিকাংশ ক্ষেত্রে !!a{derivation} পদ্ধতি ব্যবহার করি না।
বরং আমাদের অধিকাংশ প্রমাণ প্রথম-ক্রমের যুক্তিবিদ্যার ভাষায়
সম্পূর্ণ না লিখে সাধারণ ভাষায় লেখা হয় (মাঝে কিছু প্রতীকও
থাকতে পারে)। এমন প্রমাণ গড়তে প্রথমে হতবুদ্ধি লাগতে পারে—
!!a{derivation} পদ্ধতি ছাড়া কিছু প্রমাণ করব কীভাবে? শুরু করব
কোথা থেকে? প্রমাণটি ঠিক হয়েছে কি না বুঝব কীভাবে?

প্রমাণ শুরু করার আগে প্রমাণ কী এবং কীভাবে তা গড়তে হয়,
বোঝা জরুরি। নাম থেকেই বোঝা যায়, \emph{প্রমাণ} কোনো
বক্তব্য সত্য দেখানোর জন্য। কথোপকথন হিসেবে ভাবতে পারো:
কেউ জিজ্ঞাসা করল, দুই ছাড়া প্রত্যেক মৌলিক সংখ্যা বিজোড়
কি না। শুধু ``হ্যাঁ'' বললে চলবে না; \emph{কেন} তা জানতে
চাইতে পারে। তখন তুমি তাকে প্রমাণ দেবে।

দৈনন্দিন কথায় ইঙ্গিতমাত্র বা অসম্পূর্ণ উত্তরও যথেষ্ট হতে
পারে। কিন্তু যুক্তিবিদ্যা ও গণিতে চাই কঠোর প্রমাণ—যাতে
সত্যতার বিষয়ে \emph{কোনো} সংশয় না থাকে। তাই প্রমাণের
প্রত্যেক ধাপের সমর্থন দিতে হয়, এবং সেই সমর্থন যথার্থ হতে হয়
(যেমন, যে অনুমান ব্যবহার করছ, সেটি প্রমাণাধীন উপপাদ্যের
বক্তব্যে সত্যিই ধরা হয়েছে; সংজ্ঞাগুলি সঠিকভাবে প্রয়োগ
করেছ; সমর্থনের অনুমানগুলি যথার্থ, ইত্যাদি)।

সাধারণত আমরা কোনো বক্তব্য প্রমাণ করি। প্রমাণাধীন বক্তব্যের
নানা নাম আছে: প্রস্তাব, উপপাদ্য, লেমা বা অনুসিদ্ধান্ত।
প্রস্তাব হলো প্রমাণ করার মতো একটি মৌলিক বক্তব্য—লিখে রাখার
মতো গুরুত্বপূর্ণ, যদিও খুব গভীর বা বহুল ব্যবহৃত না-ও হতে
পারে। উপপাদ্য একটি তাৎপর্যপূর্ণ, গুরুত্বপূর্ণ প্রস্তাব। তার
প্রমাণ প্রায়ই কয়েক ধাপে ভাঙা হয়; কখনো প্রথম প্রমাণকারীর
নামে (যেমন কান্টরের উপপাদ্য, ল্যোভেনহাইম--স্কোলেম উপপাদ্য)
অথবা সংশ্লিষ্ট বিষয়ের নামে (যেমন পূর্ণতার উপপাদ্য) নামকরণ
হয়। লেমা হলো আরও গুরুত্বপূর্ণ কোনো ফলের প্রমাণে ব্যবহৃত
প্রস্তাব বা উপপাদ্য। বিভ্রান্তিকর হলেও, কোনো লেমা নিজেও
গুরুত্বপূর্ণ ফল হতে পারে এবং প্রবর্তকের নামে পরিচিত হতে পারে
(যেমন জর্নের লেমা)। অনুসিদ্ধান্ত হলো অন্য ফল থেকে সহজে
পাওয়া একটি ফল।

প্রমাণাধীন বক্তব্যে প্রায়ই কিছু পূর্বধারণা থাকে; এগুলি কোন
ধরনের বস্তু নিয়ে প্রমাণ করছি তা স্পষ্ট করে। বক্তব্য শুরু হতে
পারে ``$!A$ হোক $!B \lif !C$ আকারের !!a{formula}''
অথবা ``$\Gamma \Proves !A$ ধরি''—এ রকম কিছু দিয়ে।
এগুলি প্রস্তাব, উপপাদ্য বা লেমার \emph{পূর্বধারণা};
প্রমাণে সেগুলি সত্য ধরে নিতে পারো। এগুলি দাবির ক্ষেত্র সীমিত
করে এবং আলোচ্য বস্তুগুলির কিছু নামও দেয়। যেমন,
``$!A$ হোক $!B \lif !C$ আকারের !!a{formula}''
দিয়ে প্রস্তাব শুরু হলে কেবল একটি নির্দিষ্ট ধরনের
সূত্র—শর্তাধীন সূত্র—নিয়ে কিছু প্রমাণ করছ। তখন
$!B \lif !C$-কে প্রমাণে আলোচ্য যেকোনো শর্তাধীন সূত্র
হিসেবে বোঝা হয়।

\end{document}
